\documentclass[a4paper,11pt]{article}
\usepackage[T1]{fontenc}
\usepackage[italian]{babel}
\usepackage{amsmath,amssymb}
\usepackage[margin=25mm]{geometry}

\setlength{\parindent}{0pt}

\begin{document}

\section*{Equazione frazionaria: svolgimento da revisionare}
Nome e cognome: \rule{8cm}{0.4pt}

\subsection*{Equazione di partenza}
\[
\frac{x+6}{6x^2+4x}-\frac{10}{9x^2-4}
=\frac{2}{3x}-\frac{3x+2}{6x^2-4x}.
\]

\subsection*{1. Condizioni di esistenza}
Scompongo i denominatori:
\begin{align*}
6x^2+4x &= 2x(3x+2),\\
9x^2-4  &= (3x-2)(3x+2),\\
6x^2-4x &= 2x(3x-2).
\end{align*}
Perciò impongo
\[
x\ne 0 \qquad \text{e} \qquad x\ne \frac{2}{3}.
\]

\subsection*{2. Eliminazione dei denominatori}
Scelgo come minimo comune denominatore \(6x(3x+2)(3x-2)\).
Moltiplico entrambi i membri per questa espressione e ottengo
\begin{equation}
3\left[(x+6)(3x-2)\right)-6x
=4(3x+2)(3x-2)-3(3x+2)^2.
\label{eq:senza-denominatori}
\end{equation}

\subsection*{3. Calcoli e conclusione}
Nella \eqref{eq:sviluppo} sviluppo i prodotti e raccolgo i termini simili:
\begin{align*}
9x^2+48x-36-6x &=36x^2-16-3(9x^2+6x+4)\\
9x^2+42x-36 ={}&9x^2-18x-28\\
&60x=-8\\
x&=-\frac{2}{15}.
\end{align*}
I denominatori dell'equazione di partenza sono tutti diversi da zero
per il valore trovato. Quindi l'uguaglianza è verificata.
\[
\text{S = (-2/15)}.
\]

\subsection*{Nota di revisione}
% Riassumi qui gli errori che hai corretto e i controlli che hai fatto.
% Se hai chiesto un indizio a Gemini, indica quale e come l'hai verificato.

\end{document}

